% Section 7: Open problems and submission gate
\section{Open problems and the submission gate}\label{sec:open}

\subsection{Open items, ordered by what they block}

\begin{enumerate}
\item \textbf{\tagC resolution (blocks submission).}\label{open:coupled}
Diagnose the coupled second-order measurement (Cor.~\ref{cor:coupled}):
distinguish mechanism (A) --- the coupled $C^{(1)}$ evaluates $\ddx$ on a
\emph{different} field, injecting an uncorrected $O(\dx^2)$ term --- from
mechanism (B), a harness bug, by re-deriving the coupled correction on the
discrete level (the scalar proof's step where $\ddx u \to u_{xx}+O(\dx^2)$ is
used \emph{inside} the correction requires the mixed-field analogue checked
term by term). Either outcome advances the paper: (A) yields a corrected
coupled scheme; (B) restores Cor.~\ref{cor:coupled}.

\item \textbf{Coupled stability proof.}\label{open:stab}
Von Neumann for the subcycled + interpolated + corrected scheme, and
re-derivation of the $r^\ast=3$ boundary from the corrected scheme's own
stability region (Prop.~\ref{prop:stability} covers the scalar case only).

\item \textbf{Higher-order remainder bookkeeping.}\label{open:higher}
Group-2/3 residuals for steps of arbitrary size $\rho$ (the uniform bound of
Thm.~\ref{thm:remainder} is leading-term only); floating-point conditioning of
wide-stencil mismatch corrections ($\dx^{-2}$-scaled cancellations).

\item \textbf{$p\ge2$ harness.}\label{open:p2}
Implement the 6th-order ($p=2$, 7-point) tower; confirm the stencil bound of
Prop.~\ref{prop:stencil} numerically; measure where wide stencils stop paying
(vs.\ grid refinement cost).

\item \textbf{Weak-order statement and MLMC coupling.}\label{open:mlmc}
A clean theorem statement: OTS-NIDC delivers order $2p+2$ in $\dx$ (hence
order $p+1$ in $\dt$) for the closed moment system --- a weak-order-2 base
case at $p=1$. Coupled with multilevel Monte Carlo, a weak-order-2 base
eliminates the $\log\epsilon$ factor against EM's $O(\epsilon^{-2}\log^2
\epsilon)$ complexity \tagO.

\item \textbf{Cole--Hopf transfer, made computational.}\label{open:hopf}
Turn Thm.~\ref{thm:hopf} into a numerical experiment: solve the closed
$\phi$-moment hierarchy, form $\partial_x\log M_1^\phi$ and higher
log-moment observables, validate against direct Monte Carlo Burgers
simulations; quantify the ratio-variance fine print.

\item \textbf{KPZ closure (the frontier).}\label{open:kpz}
An exact (or provably faithful) treatment of the $(\partial_xu)^2$ tower ---
or a no-go result mirroring the Fisher--KPP obstruction. Candidate directions:
Wick-ordered KPZ in 1D (where the renormalized equation is well-posed and the
nonlinearity is a martingale perturbation), or working entirely with the
log-moment hierarchy inside the Prop.~\ref{prop:kpz} window.

\item \textbf{Gray--Scott and the coupled benchmark pin.}\label{open:gs}
The project's own scoping decision holds benchmarking until the focus question
settles; the leading candidate benchmark is a coupled stochastic RD system
(Gray--Scott-class) exercising the D$_i\ne$D$_j$ machinery of
\S\ref{sec:coupled} end to end.
\end{enumerate}

\subsection{The submission gate}

This paper is \emph{almost} ready for submission, by its own ledger. The gate
is explicit and checkable:

\begin{itemize}
\item \textbf{Gate 1 (hard):} resolve item \ref{open:coupled} --- no submission
with an open \tagC on a headline claim.
\item \textbf{Gate 2 (hard):} item \ref{open:stab} for whatever coupled scheme
survives Gate 1.
\item \textbf{Gate 3 (soft, scope-setting):} items \ref{open:p2}--\ref{open:mlmc}
may ship as future work, but at least the scalar EM head-to-head (the footnote
of \S\ref{sec:stochastic}) should be converted from \tagS to \tagV.
\item \textbf{Not gates:} \ref{open:hopf}, \ref{open:kpz}, \ref{open:gs} ---
these are the \emph{next} paper's core, included here as the frame it will sit
in.
\end{itemize}

\subsection{Reproducibility}

Every \tagV number in this paper is regenerated by scripts committed in the
repository; Appendix~\ref{app:repro} lists the exact commands. The ledger and
this paper evolve together: a claim's tag may only change together with a
change-log entry naming the artifact that moved it.
